\chapter{Introduction}
Probabilistic programming languages (Hakaru, Stan, Gen, etc \cite{hakaru, stan, gen}) offer a principled way of 
creating statistical models complex probability distributions
over multiple random variables, as code. The programmer (statistician) can
then interact with the model, evaluating (sampling from) the generated distribution.
This decouples the specification of a model from the inference algorithm used for sampling.
Probabilistic programming is in fact also pervasive in computer science, though generally
not cleanly abstracted into a DSL. Examples include Zero Knowledge Proof Systems (Blockchain) \cite{arklib},
Differential Privacy and Machine Learning (TensorFlow Probability \cite{tfp}). Formal reasoning about programs in
these critical applications may be highly beneficial. 

The aim of this project is to extend the reach of foundational
verification in this direction. The main contributions are:
\begin{enumerate}
\item \textbf{C1}: Implementing (formalising) Lilac \cite{lilac} and BaSL \cite{basl} separation logics in the Lean theorem prover,
with "tactical support", using Iris \cite{iris}.
\item \textbf{C2}: (if time permits) Develop a theory of compositional symbolic execution (CSE) for probabilistic programs and implement using Soteira \cite{soteira}.
CSE can automate program verification.
\end{enumerate}

% {\color{red} Here on is written hypothetically if I finish ALL of my goals}

\textbf{C1} is one of the first (if not the first) foundationally verified (in Lean) implementation of a probabilistic separation logic (PSL).
There is in fact a simultaneous project formalising the Bluebell PSL \cite{bluebell-impl} for the purpose of verifying the core of a Zero Knowledge Proof
system \cite{arklib}.

\textbf{C2} would be a novel theoretical contribution. The main
innovation here, would be either
The separation logic chosen
for this may be Lilac or BaSL, but also possibly Bluebell \cite{bluebell}, to directly apply the work in an industrial formalisation project. However
Bluebell supports both unary

with various subsystems 
But there are much wider applications
in computer science, where arguments of correctness are probabilistic in nature. Examples
expressively and formally verify these arguments if the implementation was carried out in a probabilistic
programming language/DSL, using Probabilistic Separation Logics (PSL).
This project is meant as a proof of concept, aiming to be the first implementation of a PSL.
Specifically in the Lean proof assistant, we implement the lilac separation logic
using Iris-lean (ported from rocq) to provide ergonomic tactical support. We choose
Lean for the significant body of measure-theoretic results present in Mathlib, required
to establish soundness of the logic.

Given time I would also like to pursue formalisation of BaSL or a
 theory for symbolic execution using lilac or BaSL for probabilistic programs.


on top of the Iris framework, in Lean.
offer an intuitive way to reason about
probabilistic programs at a suitable level of abstraction. 
\section{Objectives}
\section{Challenges}
\section{Contributions}